\subsubsection{アーキテクチャ記述言語とアーキテクチャ枠組み}
\keyterm{アーキテクチャ記述言語}\en{Architecture Description Language, ADL}とは、ソフトウェアアーキテクチャを表すための領域特化言語である。ADL は、もともと大規模プログラミングにおけるモジュール相互接続言語から発展した。特定の領域や様式を対象にするものもあれば、企業全体の関心事を扱う広いものもある。

UML はソフトウェア設計で広く使われるため、アーキテクチャ記述にも使われることが多い。ADL の価値は、単に図を描くことだけではない。アーキテクチャ分析、整合性確認、場合によってはコード生成を支援できる点にある。

\keyterm{アーキテクチャ枠組み}\en{architecture framework}とは、特定領域や利害関係者共同体の中で、アーキテクチャを記述するための規約、原則、実践をまとめたものである。複数のビューポイントや ADL を組み合わせ、推奨される作業方法を定める。自動車分野の AUTOSAR、統一アーキテクチャ枠組み、オープン分散処理の参照モデルなどが例である。

\begin{statementbox}{実務での使い分け}
小規模なシステムでは、軽量な図と意思決定記録で十分なことがある。一方、複数組織、規制領域、製品系列、長期運用システムでは、ADL や枠組みを使うことで整合性と説明可能性を高められる。
\end{statementbox}
